\documentclass{situationist-preprint}

\title{Brains in Vats and Semantic Externalism}
\msid{EPI-02-b}
\msfamily{EPI-02}
\msdiscipline{Epistemology}
\mstype{Companion manuscript}
\mskeywords{brain in a vat; semantic externalism; Putnam; scepticism; Moore's proof; safety}

\begin{document}
\maketitle

\begin{abstract}
Putnam argued that semantic externalism refutes the hypothesis that one is a brain in a vat, since a lifelong envatted brain could not think that it is one. EPI-02-a refutes scepticism without this argument. We show that the two are related as follows. Putnam's argument is sound for the hypothesis of lifelong envatment, but it requires premises that Moore's proof does not, and it leaves untouched the hypothesis of recent envatment, which is enough for the sceptic. Its real contribution is different: it shows that a lifelong envatted subject does not believe what we believe on the same basis, and so restricts the sceptic's counterparts to recently envatted ones. The refutation in EPI-02-a survives this restriction unchanged.
\end{abstract}

\section{Putnam's argument}

\citet[ch.~1]{Putnam1981} considered a brain that has always been in a vat, connected to a computer that supplies experiences as of an ordinary world. If reference requires a suitable causal connection to the referent, the brain's word \enquote{vat} does not refer to vats, with which it has never interacted, but to the features of the computer's program that cause its uses of the word. Hence its sentence \enquote{I am a brain in a vat} is false: it is not a brain in an image-vat. A subject who is not envatted, using \enquote{vat} to refer to vats, says something false with the same sentence. Either way the sentence is false, so I am not a brain in a vat.

\citet{Brueckner1986} gave the argument its most defensible form, using only premises about my own language. My word \enquote{vat} refers to vats; if I were a lifelong brain in a vat, my word \enquote{vat} would not refer to vats; so I am not a lifelong brain in a vat. The first premise is disquotational and known by reflection on my language; the second is semantic externalism applied to the case.

EPI-02-a does not rely on this argument (its Section~6). It argues that the subject knows, in ordinary contexts, that she has hands, because her perceptual belief is justified and safe, and that she knows she is not envatted by deduction from it (its Theorem~3.2). We ask what Putnam's argument adds.

\section{Three comparisons}

\paragraph{Premises.} Putnam's argument requires semantic externalism and knowledge of it. Externalism is a substantive and contested thesis about reference, established, if at all, by reflection on cases about which theorists disagree. Moore's proof requires only perceptual knowledge that one has hands, and the account in EPI-02-a requires only that the perceptual belief be justified and safe. A refutation of scepticism that depends on a contested semantics is less secure than one that depends on perception.

\paragraph{Scope.} The argument excludes lifelong envatment only. A subject envatted last night retains the reference her words acquired in an ordinary environment, so her sentence \enquote{I am a brain in a vat} is true, and the second premise of Brueckner's argument fails for her. The sceptic needs no more than this: the hypothesis that I was envatted last night and fed experiences continuous with my past is incompatible with my having hands, and I would have the same perceptual basis if it were true. The argument from closure, with this hypothesis in place of the lifelong one, is as strong as the original.

\paragraph{Content.} The argument nonetheless establishes something that bears on the framework of EPI-02-a. That framework assumes, for every world and basis, an envatted counterpart at which the subject believes on that basis everything she believes on it at the original world (its Assumption~2.7). If externalism is true, a lifelong envatted subject does not believe that she has hands: her word \enquote{hand} refers to features of the program, and the belief she has on her perceptual basis is a different, and true, belief. Lifelong envatted worlds are therefore not counterparts in the required sense.

\begin{proposition}\label{p:counterpart}
Given semantic externalism, Assumption~2.7 of EPI-02-a holds when the envatted counterparts are recently envatted, and fails when they are envatted for life.
\end{proposition}

\begin{proof}
A recently envatted subject's concepts are those she acquired before envatment, so on the same perceptual basis she believes the same propositions, which are false at her world. A lifelong envatted subject's concepts are individuated by her causal history with the program, so the propositions she believes on that basis differ from ours.
\end{proof}

Among the results of EPI-02-a, Assumption~2.7 is used, through Assumption~2.8, only in Theorem~3.4, which says that in the sceptic's context the subject knows neither that she has hands nor that she is not envatted. That theorem holds with recently envatted counterparts, since all it needs is a world, close in the sceptic's context, at which she believes on the same basis the very proposition that she is not envatted, and that proposition is false. The sceptic's context is therefore one in which recent envatment is salient, and Putnam's argument, by excluding lifelong envatment, merely tells the sceptic which hypothesis to attend to.

\section{Objection}

\begin{objection}
Recent envatment is a weaker sceptical hypothesis than lifelong envatment, since it leaves most of my past beliefs true. Putnam has therefore done most of the anti-sceptical work.
\end{objection}

\begin{reply}
The sceptical argument concerns my present belief that I have hands, and recent envatment falsifies it as thoroughly as lifelong envatment would. The hypothesis also iterates: whatever I now believe about the world on a perceptual basis acquired since last night is false under it. What it spares is my knowledge of the past, and scepticism about the present world is not weakened by sparing that. The anti-sceptical work is done by the safety of perceptual belief in ordinary contexts, which excludes recent envatment and lifelong envatment alike.
\end{reply}

\section{Conclusion}

Semantic externalism refutes a sceptical hypothesis, but not the one the sceptic needs. Its lasting contribution is to fix the form of the sceptical hypothesis: the counterparts that make perceptual belief unsafe in the sceptic's context are recently envatted. Moore's proof, as defended in EPI-02-a, refutes scepticism whichever hypothesis the sceptic chooses.

\begin{thebibliography}{9}\small
\bibitem[Brueckner(1986)]{Brueckner1986} Brueckner, A. (1986). Brains in a vat. \emph{The Journal of Philosophy}, 83(3), 148--167.
\bibitem[Putnam(1981)]{Putnam1981} Putnam, H. (1981). \emph{Reason, Truth and History}. Cambridge University Press.
\end{thebibliography}

\end{document}
